% Part: normal-modal-logic
% Chapter: filtrations
% Section: preliminaries

\documentclass[../../../include/open-logic-section]{subfiles}

\begin{document}

\olfileid{nml}{fil}{pre}

\olsection{પ્રારંભિક બાબતો}

ફિલ્ટ્રેશનો બતાવે છે કે આપણાં મોડલ તર્કનાં તંત્રોમાં \emph{સાન્ત
નિદર્શન ગુણધર્મ} છે, અને એ વડે તેમની નિર્ણેયતા સ્થાપિત કરી શકાય
છે. તેનો અર્થ એ કે કોઈ નિદર્શનમાં સત્ય (અસત્ય) હોય તેવું કોઈપણ
!!{formula} સાન્ત નિદર્શનમાં પણ સત્ય (અસત્ય) હોય છે.
ફિલ્ટ્રેશનો !!{formula}sના એવા ગણોને અનુલક્ષીને વ્યાખ્યાયિત થાય છે
જે ઉપસૂત્રો હેઠળ બંધ હોય.

\begin{defn}\ollabel{defn:modallyclosed}
  !!{formula}sનો ગણ $\Gamma$ \emph{ઉપસૂત્રો હેઠળ બંધ} છે જો તે
  $\Gamma$માંના દરેક !!a{formula}નાં બધાં ઉપસૂત્રો ધરાવે.
  વધુમાં, $\Gamma$ \emph{મોડલ રીતે બંધ} છે જો તે ઉપસૂત્રો હેઠળ
  બંધ હોય અને $!A \in \Gamma$ હોય ત્યારે $\Box!A,
  \Diamond!A \in \Gamma$ પણ હોય.
\end{defn}

દાખલા તરીકે, આપેલું !!a{formula}~$!A$ હોય ત્યારે તેનાં બધાં
ઉપ-!!{formula}sનો ગણ ઉપ-!!{formula}s હેઠળ બંધ છે. $!A$નાં આ
ઉપ-!!{formula}sના ગણ દ્વારા કોઈ નિદર્શનનું ફિલ્ટ્રેશન વ્યાખ્યાયિત કરીએ ત્યારે
આપણને જોઈતો ગુણધર્મ તેને હશે: મૂળ નિદર્શન $!A$ને સત્ય (અસત્ય)
બનાવે જો અને માત્ર જો તે પણ એવું કરે.

~$\Gamma$ દ્વારા~$\mModel{M}$ના ફિલ્ટ્રેશનનાં જગતોનો ગણ નીચેના
સમકક્ષતા સંબંધના બધા સમકક્ષતા-વર્ગો તરીકે વ્યાખ્યાયિત થાય છે.

\begin{defn}
ધારો કે $\mModel{M} =\tuple{W, R, V}$ અને $\Gamma$ ઉપ-!!{formula}s
હેઠળ બંધ છે. $W$ પર એવો સંબંધ $\equiv$ વ્યાખ્યાયિત કરો કે
$\Gamma$માંથી સરખાં !!{formula}sને સત્ય બનાવતાં કોઈપણ બે જગતો
તેમાં સંબંધિત હોય, એટલે કે:
\[
u \equiv v \quad \text{જો અને માત્ર જો } \quad
\forall !A \in \Gamma : \mSat{M}{!A}[u] \Leftrightarrow \mSat{M}{!A}[v].
\]
જગત~$w$નો સમકક્ષતા-વર્ગ~$[w]_\equiv$, અથવા સંક્ષેપમાં $[w]$,
એ~$w$ને $\equiv$-સમકક્ષ બધાં જગતોનો ગણ છે:
\[
[w] = \Setabs{v}{v \equiv w}.
\]
\end{defn}

\begin{prop}
  $\mModel{M}$ અને $\Gamma$ આપેલાં હોય, તો ઉપર વ્યાખ્યાયિત
  $\equiv$ સમકક્ષતા સંબંધ છે, એટલે કે તે સ્વપ્રતિબિંબિત, સંમિત
  અને પરંપરિત છે.
\end{prop}

\begin{proof}
  સંબંધ $\equiv$ સ્વપ્રતિબિંબિત છે, કારણ કે $w$ પોતે જે
  !!{formula}sને~$\Gamma$માંથી સત્ય બનાવે છે, એ જ પોતે બનાવે છે.
  તે સંમિત છે: $u$ અને $v$~$\Gamma$માંથી સરખાં !!{formula}sને
  સત્ય બનાવે, તો $v$ અને~$u$ માટે પણ એ જ સાચું છે.
  તે પરંપરિત પણ છે: $u$ અને~$v$ તથા $v$ અને~$w$ જો~$\Gamma$માંથી
  સરખાં !!{formula}sને સત્ય બનાવે, તો $u$ અને~$w$ પણ~$\Gamma$માંથી
  સરખાં !!{formula}sને સત્ય બનાવે.
\end{proof}

દરેક સમકક્ષતા સંબંધની જેમ $\equiv$ પણ~$W$ને \emph{વિભાજન-વર્ગો}માં
વહેંચે છે: એ~$W$ના એવા ઉપગણો છે જે જોડીજોડીમાં અસંયુક્ત છે અને
સાથે મળીને આખા~$W$ને આવરે છે. દરેક $w \in W$ આ વર્ગોમાંના
એકનો !!a{element} છે, એટલે કે~$[w]$નો, કારણ કે $w \equiv w$.
આમ વર્ગો $[w]$ આખા~$W$ને આવરે છે. તેઓ જોડીજોડીમાં અસંયુક્ત છે:
જો $u \in [w]$ અને $u \in [v]$ હોય, તો $u \equiv w$ અને
$u \equiv v$; તેથી સંમિતિ અને પરંપરિતતાથી $w \equiv v$,
એટલે $[w] = [v]$.

\end{document}
